% GENERATED by scripts/generate-mindmaps.mjs — do not hand-edit.
\begin{tikzpicture}[
  mindmap,
  every node/.style={
    concept,
    execute at begin node=\hskip0pt,
    align=center,
    font=\sffamily
  },
  mm root/.style={
    concept color=graphite,
    text=ghostwhite,
    font=\sffamily\bfseries,
    minimum size=3.35cm,
    text width=2.75cm
  },
  mm section/.style={
    font=\sffamily\bfseries\scriptsize,
    minimum size=2.42cm,
    text width=1.92cm
  },
  mm topic/.style={
    font=\sffamily\tiny,
    minimum size=1.92cm,
    text width=1.46cm
  }
]
  \node[mm root] (root) at (0,0) {Chapter 14\\Learning Regular Languages};
  \node[mm section, concept color=brightsky] (s1) at (54:5.1cm) {The Myhill--Nerode Foundation};
  \node[mm topic, concept color=brightsky!44!white] (t1_1) at (78:9.8cm) {The equivalence relation};
  \node[mm topic, concept color=brightsky!44!white] (t1_2) at (54:9.8cm) {Characterizing regular languages};
  \node[mm topic, concept color=brightsky!44!white] (t1_3) at (30:9.8cm) {Examples of equivalence classes};
  \node[mm section, concept color=vvteal] (s2) at (-6:5.1cm) {Passive vs. Active Learning};
  \node[mm topic, concept color=vvteal!44!white] (t2_1) at (6:9.8cm) {Passive learning and its limitations};
  \node[mm topic, concept color=vvteal!44!white] (t2_2) at (-18:9.8cm) {Active learning: The MAT framework};
  \node[mm section, concept color=vvviolet] (s3) at (-78:5.1cm) {Angluin's L* Algorithm};
  \node[mm topic, concept color=vvviolet!44!white] (t3_1) at (-42:9.8cm) {The observation table structure};
  \node[mm topic, concept color=vvviolet!44!white] (t3_2) at (-66:9.8cm) {Table properties: Minimality and Completeness};
  \node[mm topic, concept color=vvviolet!44!white] (t3_3) at (-90:9.8cm) {DFA Construction};
  \node[mm topic, concept color=vvviolet!44!white] (t3_4) at (-114:9.8cm) {The L* loop};
  \node[mm section, concept color=vvorange] (s4) at (-174:5.1cm) {Walkthrough with a Concrete Example};
  \node[mm topic, concept color=vvorange!44!white] (t4_1) at (-138:9.8cm) {Iteration 1};
  \node[mm topic, concept color=vvorange!44!white] (t4_2) at (-162:9.8cm) {Iteration 2};
  \node[mm topic, concept color=vvorange!44!white] (t4_3) at (-186:9.8cm) {Iteration 3};
  \node[mm topic, concept color=vvorange!44!white] (t4_4) at (-210:9.8cm) {Iteration 4};
  \node[mm section, concept color=tomato] (s5) at (-246:5.1cm) {Correctness, Termination, and Complexity};
  \node[mm topic, concept color=tomato!44!white] (t5_1) at (-234:9.8cm) {Complexity bounds};
  \node[mm topic, concept color=tomato!44!white] (t5_2) at (-258:9.8cm) {Proof of correctness and termination};

  \begin{scope}[on background layer, line cap=round]
    \draw[brightsky, line width=11pt] (root) -- (s1);
    \draw[brightsky!60!white, line width=5.5pt] (s1) -- (t1_1);
    \draw[brightsky!60!white, line width=5.5pt] (s1) -- (t1_2);
    \draw[brightsky!60!white, line width=5.5pt] (s1) -- (t1_3);
    \draw[vvteal, line width=11pt] (root) -- (s2);
    \draw[vvteal!60!white, line width=5.5pt] (s2) -- (t2_1);
    \draw[vvteal!60!white, line width=5.5pt] (s2) -- (t2_2);
    \draw[vvviolet, line width=11pt] (root) -- (s3);
    \draw[vvviolet!60!white, line width=5.5pt] (s3) -- (t3_1);
    \draw[vvviolet!60!white, line width=5.5pt] (s3) -- (t3_2);
    \draw[vvviolet!60!white, line width=5.5pt] (s3) -- (t3_3);
    \draw[vvviolet!60!white, line width=5.5pt] (s3) -- (t3_4);
    \draw[vvorange, line width=11pt] (root) -- (s4);
    \draw[vvorange!60!white, line width=5.5pt] (s4) -- (t4_1);
    \draw[vvorange!60!white, line width=5.5pt] (s4) -- (t4_2);
    \draw[vvorange!60!white, line width=5.5pt] (s4) -- (t4_3);
    \draw[vvorange!60!white, line width=5.5pt] (s4) -- (t4_4);
    \draw[tomato, line width=11pt] (root) -- (s5);
    \draw[tomato!60!white, line width=5.5pt] (s5) -- (t5_1);
    \draw[tomato!60!white, line width=5.5pt] (s5) -- (t5_2);
  \end{scope}
\end{tikzpicture}
